\documentclass[a4paper,12pt]{article}

\usepackage[a4paper,margin=16mm]{geometry}
\usepackage[german,main=french]{babel}
\usepackage{fontspec}
\usepackage{microtype}

\setmainfont{EBGaramond}
[
  Extension      = .otf,
  UprightFont    = *-Regular,
  ItalicFont     = *-Italic,
  BoldFont       = *-Bold,
  BoldItalicFont = *-BoldItalic,
  Numbers        = Lowercase
]

\usepackage{paracol}
\newcommand\chunks[2]{%
\begin{leftcolumn*}\begin{otherlanguage}{german}%
{#1}%
\end{otherlanguage}\end{leftcolumn*}%
\begin{rightcolumn}\begin{otherlanguage}{french}%
{#2}%
\end{otherlanguage}\end{rightcolumn}%
}

\newcommand\strophe[1]{\subsection*{\normalsize\normalfont\itshape #1}}

\begin{document}\thispagestyle{empty}
\setlength{\columnsep}{8mm}


\begin{paracol}{2}
  \chunks{\section*{Dies Bildnis ist bezaubernd schön}}
         {\section*{Ce portrait est beau à ravir}}

  \chunks{\noindent\raggedright\small Wolfgang Amadeus Mozart, \textit{Die Zauberflöte}}
         {\noindent\raggedright\small Wolfgang Amadeus Mozart, \textit{La Flûte enchantée}}

  \chunks{\noindent\raggedright\small Libretto von Emanuel Schikaneder (1791)}
         {\noindent\raggedright\small Livret d’Emanuel Schikaneder (1791)}

  \chunks{\strophe{Erster Aufzug, vierter Auftritt, Nr.~3, Arie}}
         {\strophe{Premier acte, quatrième scène, n\textsuperscript{o}~3, air}}

  \chunks{\noindent\small(Tamino ist gleich beim Empfang des Bildnisses aufmerksam geworden; seine Liebe nimmt zu, ob er gleich für alle diese Reden taub schien.)}
         {\noindent\small(Dès qu’il a reçu le portrait, Tamino l’a contemplé attentivement ; son amour grandit, bien qu’il ait paru sourd à tous ces discours.)}

  \chunks{\begin{verse}
          \textsc{Tamino}\\[\medskipamount]
          Dies Bildnis ist bezaubernd schön,\\
          Wie noch kein Auge je gesehn!\\
          Ich fühl’ es, wie dies Götterbild\\
          Mein Herz mit neuer Regung füllt.
          \end{verse}}
         {\begin{verse}
          \textsc{Tamino}\\[\medskipamount]
          Ce portrait est beau à ravir,\\
          Comme aucun œil n’en vit jamais !\\
          Je sens comme cette image divine\\
          Emplit mon cœur d’un émoi nouveau.
          \end{verse}}

  \chunks{\begin{verse}
          Dies Etwas kann ich zwar nicht nennen,\\
          Doch fühl’ ich’s hier wie Feuer brennen;\\
          Soll die Empfindung Liebe sein?\\
          Ja, ja! die Liebe ist’s allein.
          \end{verse}}
         {\begin{verse}
          Ce je-ne-sais-quoi, je ne puis le nommer,\\
          Mais je le sens ici brûler comme un feu ;\\
          Ce sentiment serait-il l’amour ?\\
          Oui, oui ! c’est l’amour, lui seul.
          \end{verse}}

  \chunks{\begin{verse}
          O wenn ich sie nur finden könnte!\\
          O wenn sie doch schon vor mir stände!\\
          Ich würde\,\dots{} würde\,\dots{} warm und rein\,\dots{}\\
          Was würde ich?\\
          Ich würde sie, voll Entzücken,\\
          An diesen heißen Busen drücken,\\
          Und ewig wäre sie dann mein.
          \end{verse}}
         {\begin{verse}
          Oh, si seulement je pouvais la trouver !\\
          Oh, si elle se tenait déjà devant moi !\\
          Je\,\dots{} je\,\dots{} avec ardeur et pureté\,\dots{}\\
          Que ferais-je ?\\
          Je la presserais, ravi,\\
          Contre ma poitrine brûlante,\\
          Et elle serait alors à moi pour toujours.
          \end{verse}}
\end{paracol}
\end{document}
